% 草案：既存の例題と同じ「例題・方針・解答・解法」の形式で検討する.
  \subsubsection{例題（草案）：整数問題とのつながり}\label{節：ガウス型の例題草案}
    ここまでに扱った多くの漸化式では,差や比を取ったり,別の数列に置き換えたりすることで,既知の形に帰着させてきた.
    ガウス記号を含む漸化式でも,この方針は重要である.
    ただし,その変形の手掛かりを得るために,整数問題で用いる考え方が必要になることがある.
    床関数の定義をもう一度書くと,
    \[
      \lfloor x\rfloor=k
      \quad\Longleftrightarrow\quad
      k\le x<k+1\quad(k\in\mathbb Z)
    \]
    である.値を求めることは,ある実数がどの二つの連続する整数の間にあるかを調べることでもある.
    また,整数との和ではガウス記号を外せたり,余りによって取りうる値を絞れたりする.
    したがって,ガウス型には共通する一つの解法があるわけではない.
    \textbf{ガウス記号が式の中で何をしているか}を読み,整数性,不等式,余り,偶奇などを利用しよう.
    一般項に加え,値の範囲や周期性を調べることにも意味がある.

    以下では,群数列として数える問題,余りを調べる問題,整数の間に挟む問題,添字を分割する問題を順に扱う.
    最後の問題はガウス記号が添字に作用する発展例であり,数列の値に作用する前の三問と区別して読もう.

    \noindent \bookblocklink{example-gauss-integer-draft}{problem}{answer}{【例題】}
    \begin{enumerate}[label=(\arabic*),parsep=3pt,format=\bookitemlink{example-gauss-integer-draft}{problem}{answer}]
      \item 数列を次のように定める.一般項を求めよ.
            \[
              a_0=1,\qquad
              a_{n+1}=a_n+\lfloor\sqrt{n+1}\rfloor+1
              \quad(n\ge0).
            \]
      \item 数列を次のように定める.一般項を求め,第2項以降の最小の正の周期を求めよ.
            \[
              a_1=2,\qquad
              a_{n+1}=X_n-3\left\lfloor\frac{X_n}{3}\right\rfloor,
            \]
            \[
              X_n=a_n^2+n^2+n\quad(n\ge1).
            \]
      \item 数列を次のように定める.ガウス記号を含まない隣接3項間漸化式を導き,一般項を求めよ.
            \[
              a_1=2,\qquad
              a_{n+1}=\lfloor(1+\sqrt3)a_n\rfloor
              \quad(n\ge1).
            \]
      \item 実数$c$に対し,数列を次のように定める.一般項を求めよ.
            \[
              a_1=1,\qquad
              a_n=a_{\left\lfloor\frac{n+1}{2}\right\rfloor}
                  +a_{\left\lfloor\frac n2\right\rfloor}+cn
              \quad(n\ge2).
            \]
            必要ならば,\,$2^m\le n<2^{m+1}$を満たす整数$m\ge0$を用いてよい.
    \end{enumerate}

    \noindent \bookblocklink{example-gauss-integer-draft}{hint}{problem}{【方針】}
    \begin{enumerate}[label=(\arabic*),parsep=3pt,format=\bookitemlink{example-gauss-integer-draft}{hint}{problem}]
      \item \sectionref{節：階差型}と同じように階差を足します.
            \,$\lfloor\sqrt{k}\rfloor=j$となる範囲を平方数で区切り,各群の項数を数えましょう.
      \item \,$X_n$が整数であることを確かめ,\,$X_n-3\lfloor\frac{X_n}{3}\rfloor$を3で割った余りとして読みます.
            値だけでなく,添字を3で割った余りも一緒に追いましょう.
      \item \,$\lambda=1+\sqrt3$とおき,丸め誤差$\varepsilon_n=\lambda a_n-a_{n+1}$を調べます.
            \,$\lambda^2=2\lambda+2$を使うと,次の項を求める式が「整数から1未満の正の数を引く形」になります.
      \item 添字の偶奇で分けてから,階差$b_n=a_{n+1}-a_n$を取ります.
            \,$b_{2k}$と$b_{2k+1}$を$b_k$で表し,\,$2^m\le n<2^{m+1}$の範囲を一つの群として考えましょう.
    \end{enumerate}

    \noindent \bookblocklink{example-gauss-integer-draft}{answer}{problem}{【解答】}
    \begin{enumerate}[label=(\arabic*),parsep=3pt,format=\bookitemlink{example-gauss-integer-draft}{answer}{problem}]
      \item \,$m=\lfloor\sqrt n\rfloor$とおくと,
            \[
              a_n=1+n+m(n+1)
                  -\frac{m(m+1)(2m+1)}6
              \quad(n\ge0).
            \]
      \item 初項は$a_1=2$であり,\,$k\ge0$について,
            \[
              \begin{aligned}
                a_{6k+2}=a_{6k+3}=a_{6k+4}&=0,\\
                a_{6k+5}&=2,\\
                a_{6k+6}=a_{6k+7}&=1.
              \end{aligned}
            \]
            第2項以降の最小の正の周期は6である.
      \item \,$a_1=2,\ a_2=5$であり,
            \[
              a_{n+2}=2a_{n+1}+2a_n-1.
            \]
            \,$\alpha=1+\sqrt3,\ \beta=1-\sqrt3$とおくと,
            \[
              a_n=\frac13
                +\frac{5+3\sqrt3}{6}\alpha^{n-1}
                +\frac{5-3\sqrt3}{6}\beta^{n-1}.
            \]
      \item \,$2^m\le n<2^{m+1}$を満たす整数$m\ge0$を用いると,
            \[
              a_n=n+c\{(m+2)n-2^{m+1}\}
              \quad(n\ge1).
            \]
    \end{enumerate}

    \noindent \bookblocklink{example-gauss-integer-draft}{solution}{problem}{【解法】}
    \begin{enumerate}[label=(\arabic*),parsep=3pt,format=\bookitemlink{example-gauss-integer-draft}{solution}{problem}]
      \item \textbf{同じ値を取る添字の範囲を調べる.}\\
            階差を足すと,
            \[
              a_n=1+n+\sum_{k=1}^{n}\lfloor\sqrt k\rfloor
              \quad(n\ge0)
            \]
            となる.ここで$n=0$の和は0とする.
            \,$m=\lfloor\sqrt n\rfloor$とおくと,
            \[
              m^2\le n<(m+1)^2.
            \]
            \,$\lfloor\sqrt k\rfloor=j$となる条件は,
            \[
              j^2\le k\le(j+1)^2-1
            \]
            であるから,第$j$群には$2j+1$項がある.
            \,$m\ge1$なら第$m-1$群まではすべて足し,第$m$群では$k=m^2$から$n$までの$n-m^2+1$項を足せばよい.
            したがって,
            \begin{align*}
              \sum_{k=1}^{n}\lfloor\sqrt k\rfloor
                &=\sum_{j=1}^{m-1}j(2j+1)+m(n-m^2+1)\\
                &=\frac{(m-1)m(4m+1)}6+m(n-m^2+1)\\
                &=m(n+1)-\frac{m(m+1)(2m+1)}6.
            \end{align*}
            これを代入すると,【解答】の一般項を得る.\,$n=0$では$m=0$であり,同じ式が$a_0=1$を与える.
            群の境界を不等式で求め,整数の個数を数えることが,この問題の整数問題的な着眼点である.

            \noindent \textbf{【別解】}\\
            \,$\lfloor\sqrt k\rfloor$は,\,$j^2\le k$を満たす正整数$j$の個数でもある.
            \,$j=1,\ldots,m$を固定すると,\,$j^2\le k\le n$となる$k$は$n-j^2+1$個ある.
            同じ組$(j,k)$を$j$ごとに数え直すと,
            \begin{align*}
              \sum_{k=1}^{n}\lfloor\sqrt k\rfloor
                &=\sum_{j=1}^{m}(n-j^2+1)\\
                &=m(n+1)-\frac{m(m+1)(2m+1)}6.
            \end{align*}
            同じ値の群をまとめる方法と,各項に含まれる1を数え直す方法の二通りがある.

      \item \textbf{余りとして読み,値と添字の余りを組にして追う.}\\
            初項は整数であり,整数$a_n$から作られる$X_n$も整数である.
            さらに,\,$X_n-3\lfloor\frac{X_n}{3}\rfloor$は整数なので,帰納法によりすべての項が整数である.
            床関数の定義から,
            \[
              0\le X_n-3\left\lfloor\frac{X_n}{3}\right\rfloor<3
            \]
            である.よって第2項以降は$a_n\in\{0,1,2\}$となり,
            \[
              a_{n+1}\equiv a_n^2+n(n+1)\pmod3
            \]
            によって,次の値が一意に決まる.
            ここで,合同式$x\equiv y\pmod3$は,\,$x-y$が3の倍数であることを表す.
            また,\,$n(n+1)$を3で割った余りは,\,$n$の余りが0,1,2のとき,それぞれ0,2,0である.
            実際に追うと,
            \[
              a_1,a_2,\ldots,a_8=2,0,0,0,2,1,1,0.
            \]
            第2項と第8項では値がともに0であり,添字も3で割った余りが2で一致する.
            同じ値と同じ添字の余りからは同じ次の値が決まるので,帰納法により,
            \[
              a_{n+6}=a_n\quad(n\ge2)
            \]
            となる.したがって,【解答】の一般項が得られる.
            値$a_n$だけの一致では十分ではない.次の値には$n$も関わるので,添字の余りを合わせて追う必要がある.

            最小周期も確かめよう.第2項以降では,\,$a_2=0$から次に値2が現れるのは$a_5$であり,その後は6項ごとに現れる.
            正の周期$p$があるなら,\,$a_{5+p}=a_5=2$であるから$p$は6の倍数である.
            6は実際に周期なので,最小の正の周期は6である.
            なお,\,$a_1=2\ne a_7=1$なので,初項から6周期となるわけではない.

      \item \textbf{丸め誤差を評価し,ガウス記号の値を確定する.}\\
            \,$\lambda=1+\sqrt3$とおく.初項は正の整数であり,正の整数$a_n$に対して$\lambda a_n>2$なので,次の項も正の整数である.
            また,\,$\lambda$は無理数なので$\lambda a_n$は整数ではない.
            そこで,
            \[
              \varepsilon_n=\lambda a_n-a_{n+1}
            \]
            とおくと,床関数の定義から$0<\varepsilon_n<1$である.
            この範囲を使って,次に床を取る数がどの整数の間にあるかを調べる.
            \,$\lambda^2=2\lambda+2$より,
            \begin{align*}
              \lambda a_{n+1}
                &=\lambda(\lambda a_n-\varepsilon_n)\\
                &=(2\lambda+2)a_n-\lambda\varepsilon_n\\
                &=2a_{n+1}+2a_n-(\lambda-2)\varepsilon_n.
            \end{align*}
            \,$0<\lambda-2=\sqrt3-1<1$なので,
            \[
              0<(\lambda-2)\varepsilon_n<1.
            \]
            整数$M_n=2a_{n+1}+2a_n$を用いれば,
            \[
              M_n-1<\lambda a_{n+1}<M_n
            \]
            である.ここで床を取ると,
            \[
              a_{n+2}=M_n-1=2a_{n+1}+2a_n-1.
            \]
            初期条件は$a_1=2,\ a_2=\lfloor2+2\sqrt3\rfloor=5$である.
            丸め誤差そのものを求めなくても,その範囲から床関数の値を確定できる.

            あとは既知の形に帰着させる.定数項を消すため,
            \[
              b_n=a_n-\frac13
            \]
            とおくと,
            \[
              b_{n+2}=2b_{n+1}+2b_n,
              \quad b_1=\frac53,\quad b_2=\frac{14}3.
            \]
            \sectionref{節：隣接3項間漸化式型}で学んだように,特性方程式$x^2-2x-2=0$の2解は,
            \[
              \alpha=1+\sqrt3,\qquad\beta=1-\sqrt3
            \]
            であるから,\,$b_n=A\alpha^{n-1}+B\beta^{n-1}$と書ける.
            初期条件から,
            \[
              A+B=\frac53,\qquad
              \alpha A+\beta B=\frac{14}3
            \]
            を解くと,
            \[
              A=\frac{5+3\sqrt3}{6},\qquad
              B=\frac{5-3\sqrt3}{6}.
            \]
            よって,【解答】の一般項を得る.

      \item \textbf{添字の偶奇で分け,階差を2の累乗ごとにまとめる.}\\
            \,$n\ge2$では床関数が指定する二つの添字はともに1以上$n-1$以下なので,数列は順に定義できる.
            偶数と奇数に分けると,\,$k\ge1$について,
            \[
              \begin{aligned}
                a_{2k}&=2a_k+2ck,\\
                a_{2k+1}&=a_{k+1}+a_k+c(2k+1).
              \end{aligned}
            \]
            両式を引けば,階差が現れる.\,$b_n=a_{n+1}-a_n$とおくと,
            \[
              b_{2k}=a_{2k+1}-a_{2k}=b_k+c.
            \]
            また,\,$a_{2k+2}=2a_{k+1}+2c(k+1)$より,
            \[
              b_{2k+1}=a_{2k+2}-a_{2k+1}=b_k+c.
            \]
            初期値は$a_2=2+2c$から$b_1=1+2c$である.
            したがって,
            \[
              b_{2k}=b_{2k+1}=b_k+c
            \]
            となり,階差の添字も半分にできる.
            第$m$群を$2^m\le n<2^{m+1}$の範囲とすると,
            \[
              b_n=1+(m+2)c
            \]
            である.これを$m$について帰納法で確かめよう.
            \,$m=0$では$n=1$のみであり,\,$b_1=1+2c$と一致する.
            第$m$群の値が分かっていれば,第$m+1$群の添字は$2k$または$2k+1$と書け,その$k$は第$m$群に属する.
            よって,次の群の値は$c$だけ増え,\,$1+(m+3)c$となる.

            第$j$群には$2^j$項ある.\,$2^m\le n<2^{m+1}$のとき,\,$b_1$から$b_{n-1}$までには第$m-1$群までの全項と,第$m$群の$n-2^m$項が含まれる.
            従って,
            \begin{align*}
              a_n
                &=1+\sum_{j=0}^{m-1}2^j\{1+(j+2)c\}\\
                &\qquad +(n-2^m)\{1+(m+2)c\}.
            \end{align*}
            ここで,
            \[
              \sum_{j=0}^{m-1}2^j=2^m-1,\qquad
              \sum_{j=0}^{m-1}j2^j=(m-2)2^m+2
            \]
            を使う.後者は,和を$T$とおいて$2T-T$を計算すれば得られる.
            \,$m=0$ではいずれも空の和として0であり,右辺も0になる.
            代入して整理すると,
            \[
              a_n=n+c\{(m+2)n-2^{m+1}\}
            \]
            を得る.添字を分割する問題も,階差を取ることで群数列の総和につながった.
    \end{enumerate}
